% Section 6: the bounded dual simplex.
\section{The bounded dual simplex method}\label{sec:simplex}

The dual simplex is \nirnay{}'s default LP engine and the node solver of branch-and-bound. Its
design follows Koberstein's thesis~\cite{koberstein2005}: dual steepest-edge pricing
\cite{forrest1992}, the bound-flipping ratio test \cite{maros2003} with Harris's two-pass
tolerance \cite{harris1973}, Fourer's auxiliary problem for phase~1 \cite{fourer1994}, cost
perturbation and shifting against degeneracy, and a primal simplex to clean up at the end
(\file{lp/simplex.py}, with the per-iteration loops in \file{lp/\_simplex\_kernels.py}).

\subsection{Basis, primal and dual values}
In the computational form \eqref{eq:simplex-form} the constraint matrix is $\bar A=[A\;\;{-I}]$ over
$n+m$ variables. A basis is an index set $\mathcal B$ of size $m$ with $B=\bar A_{:,\mathcal B}$
nonsingular; every nonbasic variable sits at a bound (or at zero if free). Then
\begin{equation}\label{eq:basic}
  x_{\mathcal B} = -B^{-1}\bar A_{:,\mathcal N}\,x_{\mathcal N},\qquad
  y = B\mT c_{\mathcal B},\qquad d = c - \bar A\T y\quad(d_{\mathcal B}=0).
\end{equation}
The basis is \emph{dual feasible} when every nonbasic reduced cost has the sign its bound allows
($d_j\ge0$ at a lower bound, $d_j\le0$ at an upper bound, $d_j=0$ if free), and
\emph{primal feasible} when $\ell_{\mathcal B}\le x_{\mathcal B}\le u_{\mathcal B}$. The dual
simplex keeps dual feasibility and removes primal infeasibility one row at a time. Each iteration
does not decrease the dual objective \eqref{eq:dualobj}.

\subsection{One iteration}
\paragraph{Pricing: dual steepest edge}
Let $\delta_r$ be the primal infeasibility of basic variable $r$ ($\ell_r-x_r$ below, $x_r-u_r$
above; only violations beyond the primal tolerance $10^{-7}$ count). The leaving row maximises
\begin{equation}\label{eq:dse}
  r = \arg\max_i \frac{\delta_i^2}{w_i},\qquad w_i \approx \bigl\|e_i\T B^{-1}\bigr\|_2^2 ,
\end{equation}
which measures the infeasibility per unit length of the dual edge~\cite{forrest1992}. The weights start
at~1 for the slack basis.

\paragraph{Pivot row}
One BTRAN gives $\rho_r=B\mT e_r$. The pivot row $\alpha_r=\rho_r\T\bar A$ is then formed row-wise: the
kernel walks only the rows of $A$ (columns of $A\T$) where $\rho_r$ is nonzero, so a sparse $\rho_r$
costs little. Basic entries are set to zero.

\paragraph{Ratio test}
If $x_r$ leaves towards its lower bound, set $s=-1$; towards its upper bound, $s=+1$. A nonbasic $j$ is
a candidate when $s\,\alpha_{rj}$ has the sign that lets $d_j$ move towards zero: $s\alpha_{rj}>\epsilon$
at a lower bound, $s\alpha_{rj}<-\epsilon$ at an upper bound, $|\alpha_{rj}|>\epsilon$ if free. The
pivot tolerance is \emph{relative}, $\epsilon=\max(10^{-7},\,10^{-9}\max_j|\alpha_{rj}|)$, so a
badly scaled row cannot supply a pivot that is tiny next to its neighbours. Each candidate has a
breakpoint $t_j=|d_j|/|\alpha_{rj}|$ (negative $d_j$ of the wrong sign is clipped to zero) and a
relaxed breakpoint $\tilde t_j=(|d_j|+\epsilon_d)/|\alpha_{rj}|$ with dual tolerance
$\epsilon_d=10^{-7}$. \Cref{alg:ratio} passes the breakpoints in increasing order, as
Maros's bound-flipping ratio test does~\cite{maros2003}. The slope of the dual objective starts at
$\delta_r$ and falls by $|\alpha_{rj}|(u_j-\ell_j)$ each time a boxed candidate is passed and flipped
to its other bound. The pass stops before the slope would turn negative. Among the remaining
candidates, Harris's two passes~\cite{harris1973} choose the entering~$q$: pass~1 takes the bound
$\Theta=\min\tilde t_j$, and pass~2 takes, among the candidates with $t_j\le\Theta$, the one with the
largest $|\alpha_{rj}|$. Stepping to any such candidate leaves every reduced cost within $\epsilon_d$
of feasibility, and the large pivot keeps the basis well conditioned.

\begin{algorithm}[t]
\caption{Bound-flipping ratio test with Harris's two passes (\code{dual\_ratio\_test})}\label{alg:ratio}
\Input{pivot row $\alpha_r$, reduced costs $d$, statuses, bounds, direction $s$, infeasibility $\delta$, tolerances $\epsilon_d,\epsilon_p$}
\Output{entering column $q$ (or none), list $F$ of bound flips}
$\epsilon\leftarrow\max(\epsilon_p,\ 10^{-9}\max_j|\alpha_{rj}|)$ \Comment*[r]{relative pivot tolerance}
$\mathcal C\leftarrow$ candidates with $s\alpha_{rj}$ of the right sign and $|\alpha_{rj}|>\epsilon$\;
\lIf{$\mathcal C=\emptyset$}{\Return none (the dual is unbounded along this row)}
sort $\mathcal C$ by $t_j=|d_j|/|\alpha_{rj}|$;\quad $\text{slope}\leftarrow\delta$;\quad $F\leftarrow\emptyset$\;
\For{$j\in\mathcal C$ in order, except the last}{
  \lIf{$j$ is free or $u_j-\ell_j=\infty$ or $\text{slope}-|\alpha_{rj}|(u_j-\ell_j)\le0$}{\textbf{break}}
  $F\leftarrow F\cup\{j\}$;\quad $\text{slope}\leftarrow\text{slope}-|\alpha_{rj}|(u_j-\ell_j)$\;
}
$\mathcal R\leftarrow$ the candidates not in $F$\;
$\Theta\leftarrow\min_{j\in\mathcal R}(|d_j|+\epsilon_d)/|\alpha_{rj}|$ \Comment*[r]{Harris pass 1}
$q\leftarrow\arg\max\{|\alpha_{rj}| : j\in\mathcal R,\ t_j\le\Theta\}$ \Comment*[r]{Harris pass 2}
\Return $q$, $F$\;
\end{algorithm}

\paragraph{Updates}
FTRAN gives the entering column $\alpha_q=B^{-1}\bar a_q$. Its entry $\alpha_{q,r}$ must agree with
$\alpha_{rq}$ from the row: if $|\alpha_{q,r}-\alpha_{rq}|>10^{-6}(1+|\alpha_{q,r}|)$ or
$|\alpha_{q,r}|<10^{-9}$, the factorisation has lost accuracy. The basis is then refactorised and the
iteration repeated; if the factorisation was already fresh, row $r$ is skipped for now. Otherwise the
flips move the basic variables once, $x_{\mathcal B}\leftarrow x_{\mathcal B}-B^{-1}\sum_{j\in F}\bar a_j\Delta x_j$, and
\begin{equation}\label{eq:updates}
\begin{aligned}
  \theta_d &= d_q/\alpha_{rq}, & d_j &\leftarrow d_j-\theta_d\,\alpha_{rj}\ (j\in\mathcal N), & d_{\text{leave}}&=-\theta_d,\\
  \theta_p &= (x_{\mathcal B_r}-\text{target})/\alpha_{q,r}, & x_{\mathcal B}&\leftarrow x_{\mathcal B}-\theta_p\,\alpha_q, & x_q&\leftarrow x_q+\theta_p ,
\end{aligned}
\end{equation}
where the leaving variable goes to the bound it violated. The dual steepest-edge weights are
updated exactly as in Forrest and Goldfarb~\cite{forrest1992}, with one more FTRAN $\tau=B^{-1}\rho_r$
and $\beta_r=\|\rho_r\|^2$:
\begin{equation}\label{eq:dse-update}
  w_i\leftarrow \max\Bigl(w_i + \Bigl(\frac{\alpha_{iq}}{\alpha_{rq}}\Bigr)^{\!2}\beta_r - 2\frac{\alpha_{iq}}{\alpha_{rq}}\tau_i,\;
  10^{-4}\Bigl(1+\Bigl(\frac{\alpha_{iq}}{\alpha_{rq}}\Bigr)^{\!2}\Bigr)\Bigr)\ (i\ne r),\qquad
  w_r\leftarrow\max\Bigl(\frac{\beta_r}{\alpha_{rq}^2},10^{-8}\Bigr).
\end{equation}
The lower bounds on the weights keep rounding from driving a weight to zero or below. Finally $q$
replaces the leaving variable in the basis, and an eta vector is appended to the factorisation
(\cref{sec:lu}). \Cref{fig:simplex-flow} shows the whole loop.

\begin{figure}[p]
\centering
\begin{tikzpicture}[node distance=5.5mm and 8mm]
  \node[term] (start) {start: basis $\mathcal B$, bounds, costs};
  \node[blk,below=of start,minimum width=5.4cm] (fac) {refactorise $B$ (or reuse LU + etas)\\[-1pt]{\scriptsize recompute $x_{\mathcal B}$, $y$, $d$; repair wrong-sign $d_j$}};
  \node[blk,below=of fac,minimum width=5.4cm] (price) {pricing: $r=\arg\max\delta_i^2/w_i$ \scriptsize(DSE)};
  \node[dec,below=of price] (any) {row\\found?};
  \node[term,fill=igreen,draw=igreen,right=14mm of any] (opt) {optimal (phase done)};
  \node[blk,below=of any,minimum width=5.4cm] (row) {BTRAN $\rho_r=B\mT e_r$;\ pivot row $\alpha_r=\rho_r\T[A\;{-I}]$};
  \node[blk,below=of row,minimum width=5.4cm] (ratio) {bound-flipping ratio test + Harris\\[-1pt]{\scriptsize\cref{alg:ratio}: entering $q$, flips $F$}};
  \node[dec,below=of ratio] (cand) {$q$\\exists?};
  \node[dec,right=14mm of cand] (fresh) {fresh\\factor?};
  \node[term,fill=saffron,draw=saffron,right=8mm of fresh] (inf) {primal infeasible};
  \node[blk,below=of cand,minimum width=5.4cm] (ftran) {FTRAN $\alpha_q=B^{-1}\bar a_q$};
  \node[dec,below=of ftran] (agree) {$\alpha_{q,r}\approx\alpha_{rq}$?};
  \node[blkm,left=10mm of agree,text width=2.4cm] (skip) {refactorise;\\[-1pt]{\scriptsize skip row $r$ if factor was fresh}};
  \node[blk,below=of agree,minimum width=5.4cm] (upd) {apply flips; update $d$, $x_{\mathcal B}$, $w$\\[-1pt]{\scriptsize\eqref{eq:updates}, \eqref{eq:dse-update}}};
  \node[blk,below=of upd,minimum width=5.4cm] (chg) {basis change; append eta};
  \node[dec,below=of chg] (due) {refactor\\due?};
  \draw[flow] (start) -- (fac);
  \draw[flow] (fac) -- (price);
  \draw[flow] (price) -- (any);
  \draw[flow] (any) -- node[lbl,above]{no} (opt);
  \draw[flow] (any) -- node[lbl,right]{yes} (row);
  \draw[flow] (row) -- (ratio);
  \draw[flow] (ratio) -- (cand);
  \draw[flow] (cand) -- node[lbl,above]{no} (fresh);
  \draw[flow] (fresh) -- node[lbl,above]{yes} (inf);
  \draw[flow] (fresh.north) -- node[lbl,right,pos=0.3]{no: confirm} ++(0,1.0) |- ([yshift=-1mm]fac.east);
  \draw[flow] (cand) -- node[lbl,right]{yes} (ftran);
  \draw[flow] (ftran) -- (agree);
  \draw[flow] (agree) -- node[lbl,above]{no} (skip);
  \draw[flow] (skip.north) |- ([yshift=1mm]fac.west);
  \draw[flow] (agree) -- node[lbl,right]{yes} (upd);
  \draw[flow] (upd) -- (chg);
  \draw[flow] (chg) -- (due);
  \draw[flow] (due.east) -- node[lbl,above,pos=0.2]{yes} ++(2.6,0) |- ([yshift=2mm]fac.east);
  \draw[flow] (due.west) -- node[lbl,above,pos=0.2]{no} ++(-3.9,0) |- (price.west);
\end{tikzpicture}
\caption[Dual simplex iteration loop]{The dual simplex iteration loop of \file{lp/simplex.py}.
Diamonds are tests. Every path that could declare a result on a possibly inaccurate
factorisation (a missing ratio-test candidate, a pivot that the row and the column see
differently) first refactorises and repeats the iteration.}
\label{fig:simplex-flow}
\end{figure}

\subsection{Phase 1, perturbation and cleanup}
\paragraph{Dual phase 1 by auxiliary boxes}
The all-logical start basis is rarely dual feasible. Following Fourer~\cite{fourer1994}
(Koberstein~\cite[§4.1.2]{koberstein2005}), phase~1 solves the same problem with every bound replaced by
a box:
\begin{equation}\label{eq:aux}
  \text{free}\to[-1,1],\qquad \text{lower only}\to[0,1],\qquad \text{upper only}\to[-1,0],\qquad
  \text{boxed or fixed}\to[0,0].
\end{equation}
Every variable of the auxiliary problem is boxed, so any basis becomes dual feasible once each
nonbasic variable is placed at the bound its reduced cost prefers. The dual simplex therefore starts
at once. An optimal basis of the auxiliary problem is dual feasible for the real problem whenever the
real problem is dual feasible at all. The code does not test the auxiliary objective. Phase~2
starts from the phase-1 basis, and any reduced cost left with the wrong sign is repaired by the
flips and cost shifts described below. If the real problem is dual infeasible (unbounded), the
shifted costs are removed at the end and the primal cleanup finds the unbounded ray.

\paragraph{Cost perturbation}
Degenerate LPs make many ratio-test breakpoints tie at zero, and the dual simplex can then cycle
or stall. Both phases therefore run on perturbed costs~\cite[§6.2.2.3]{koberstein2005}:
\begin{equation}\label{eq:perturb}
  \tilde c_j = c_j + \sigma_j\,5\cdot10^{-7}(1+|c_j|)(1+\xi_j),\qquad \xi_j\sim U[0,1],
\end{equation}
with $\sigma_j=-1$ for a variable at its upper bound and $+1$ otherwise. The direction keeps the
start dual feasible, and the random magnitude breaks the ties. Basic, fixed and free nonbasic
variables are not perturbed. The random factors $\xi$ come from a generator with a fixed seed and are
drawn once per instance, so runs are reproducible and the warm-started re-solves of branch-and-bound
all see the same perturbation.

\paragraph{Cost shifting}
If the entering candidate's reduced cost has the wrong sign (by less than the tolerance, which
the relaxed test allows), its cost is shifted so that $d_q=0$ before the step. The dual objective
then never moves backwards~\cite[§6.2.2.2]{koberstein2005}. After each refactorisation the same fix
is applied to all nonbasics: a boxed variable with a wrong-sign reduced cost is flipped to its other
bound, and any other variable has its cost shifted.

\paragraph{Cleanup by primal simplex}
When phase~2 is optimal, the perturbation and shifts are removed: the true costs are restored and
$y$, $d$ recomputed. Any dual infeasibility that remains is removed by a primal simplex on the same
factorisation. It prices by the largest dual infeasibility, uses a Harris two-pass ratio test over the
basic variables, and flips the entering variable instead of pivoting when it reaches its own
other bound first. The cleanup runs with the dual tolerance tightened to $10^{-9}$
(\cref{sec:simplex-safeguards}). If it loses primal feasibility by more than $10\epsilon_p$, the dual
simplex runs once more. On return the duals are recomputed exactly as $y=B\mT c_{\mathcal B}$ for
the true costs.

\subsection{Warm starts for branch-and-bound}
A \code{SimplexLP} object keeps its factorisation, basis, weights and phase-1 status between calls.
\code{set\_structural\_bounds} changes only the variables whose bounds differ, and re-places
nonbasic ones at their new bounds. \code{get\_basis} and \code{load\_basis} save and restore a basis as
one status byte per variable. If the loaded basis has the same basic set as the current one, the
LU and eta file remain valid and only $x_{\mathcal B}$, $y$ and $d$ are recomputed; otherwise the
basis is refactorised. After a branching change the parent's optimal basis stays dual feasible, so a
child node is typically re-optimised in a few dual pivots.
